\documentclass[../../main.tex]{subfiles}
\begin{document}
\section{Second variation, Reilly identities, and the splitting mechanism}
\label{sec:jacobi}

This section assembles the geometric mechanisms intended to prove the weighted Stein-trace
estimate for near-minimal cuts. The framework is classical; the statements proved here are the
exact model statements, and the quantitative versions are posed as open targets in
Section~\ref{sec:open}. Throughout, $\nu=e^{-V}dx$ is smooth log-concave with smooth convex
support, and $\Sigma=\partial^*E$ with the notation of Section~\ref{sec:notation}.

\subsection{Second variation and stability of near-minimizers}

\begin{proposition}[Second variation; stability inequality]
\label{prop:second-variation}
Let $E$ be a smooth critical set of the weighted perimeter under the volume constraint, so
that the weighted mean curvature $H_\nu=H-\partial_nV$ is constant on $\Sigma$. For a normal
variation with speed $u$ satisfying the linearized volume constraint
$\int_\Sigma u\dd\sigma_\nu=0$, the second variation of weighted perimeter is the index form
$\calI_\Sigma(u,u)$ of \eqref{eq:jacobi-form} with Jacobi potential
$\qJac_\Sigma=\abs{II_\Sigma}^2+\nabla^2V(n,n)$. If $E$ is a local minimizer, then
\begin{equation}\label{eq:stability}
  \calI_\Sigma(u,u)\ \ge\ 0
  \qquad\text{for all }u\in C_c^\infty(\Sigma)\text{ with }\int_\Sigma u\dd\sigma_\nu=0 .
\end{equation}
For log-concave $\nu$ one has $\qJac_\Sigma\ge\abs{II_\Sigma}^2\ge0$, so the constant mode is
always destabilizing: $\calI_\Sigma(1,1)=-\mathfrak K_\Sigma\le0$, with $\mathfrak K_\Sigma$
the constant-mode curvature of \eqref{eq:constant-mode-K}. Stability of minimizers is carried
entirely by the volume constraint, which removes the constant mode.
\end{proposition}

\begin{proof}
The first and second variation formulas for weighted perimeter under a volume constraint are
classical; see \cite{Bayle2004,BayleRosales2003,Rosales2014,Milman2009Isoperimetric} for the
weighted (manifold-with-density) setting, including the free-boundary contribution
$\calB_{\partial K}$ when $\Sigma$ meets the boundary of the support. The sign assertions are
immediate from $\nabla^2V\succeq0$.
\end{proof}

\begin{remark}[The constant mode is the obstruction, again]
\label{rem:constant-mode}
The stability inequality \eqref{eq:stability} gives quantitative control of all mean-zero modes
of a near-minimal boundary: this is the engine by which stability theory should bound the
non-constant part of the boundary trace $\partial_nu_M$ in \eqref{eq:trace-CS}. The constant
mode is invisible to \eqref{eq:stability} --- it is removed by the constraint --- and its
size is governed by $\mathfrak K_\Sigma$. This is the same residual already isolated twice: as
the operator-to-trace upgrade in Section~\ref{sec:carleson} and as the projection-test ceiling
in Section~\ref{sec:qcts}. The three formulations are different coordinates on one obstruction,
and the splitting mechanism below is the geometric proposal for handling it.
\end{remark}

\subsection{Constant-mode trace control and profile curvature}

The following elementary Schur-complement estimate is the local form of the constant-mode problem.
Let
\[
  P=\sigma_\nu(\Sigma),\qquad
  T(\psi)=\int_\Sigma\psi\dd\sigma_\nu,\qquad
  \mathfrak K_\Sigma=-\calI_\Sigma(1,1)>0 .
\]
For $\psi=\bar\psi+\psi_0$ with $\bar\psi=P^{-1}T(\psi)$ and
$\int_\Sigma\psi_0\dd\sigma_\nu=0$, stability of $\psi_0$ gives the quadratic constraint
\[
  \mathfrak K_\Sigma\bar\psi^2
  +2\calI_\Sigma(\psi,1)\bar\psi-\calI_\Sigma(\psi,\psi)\le0 .
\]
Solving it yields
\begin{equation}\label{eq:Jsharp-control}
  \abs{T(\psi)}^2
  \le
  \frac{2P^2}{\mathfrak K_\Sigma}
  \left(
    \calI_\Sigma(\psi,\psi)
    +\frac{2}{\mathfrak K_\Sigma}\abs{\calI_\Sigma(\psi,1)}^2
  \right).
\end{equation}
Thus the relevant boundary energy is
\begin{equation}\label{eq:Jsharp-def}
  \Jsharp_\Sigma(\psi)
  =
  \calI_\Sigma(\psi,\psi)
  +\frac{2}{\mathfrak K_\Sigma}\abs{\calI_\Sigma(\psi,1)}^2 .
\end{equation}
The coefficient $2/\mathfrak K_\Sigma$ is not cosmetic: with coefficient
$1/\mathfrak K_\Sigma$, pure constants are not controlled. In the Gaussian halfspace model
$\mathfrak K_\Sigma=P$ and \eqref{eq:Jsharp-control} is sharp at the level of constants.

The same curvature controls the local concavity of the isoperimetric profile.

\begin{lemma}[Profile bound]
\label{lem:profile-bound}
Suppose that at volume $p$ there is a smooth volume-constrained minimizer $E_p$ with boundary
$\Sigma_p$, weighted area $P=I(p)$, and that $I$ is twice differentiable at $p$. Then
\begin{equation}\label{eq:profile-bound}
  \mathfrak K_{\Sigma_p}\le -I''(p)I(p)^2 .
\end{equation}
\end{lemma}

\begin{proof}
Flow $\Sigma_p$ with unit normal speed. Let $v(s)$ be the enclosed volume and $P(s)$ the weighted
area. Then $v'(0)=P$ and $P'(0)=\lambda P$, where $\lambda=H_\nu$ is the constant weighted mean
curvature. The Riccati variation of weighted mean curvature gives
$P''(0)=-\mathfrak K_{\Sigma_p}+\lambda^2P$, with support terms included in
$\mathfrak K_{\Sigma_p}$. Writing $\Psi(w)=P(v^{-1}(w))$ gives a smooth upper barrier for $I$ at
$p$, and the chain rule gives $\Psi''(p)=-\mathfrak K_{\Sigma_p}/P^2$. Since $I\le\Psi$ and
$I(p)=\Psi(p)$, twice differentiability gives $I''(p)\le\Psi''(p)$.
\end{proof}

\begin{corollary}[Constant-mode degeneracy in the small-Cheeger regime]
\label{cor:generic-degeneracy}
Grant smooth minimizers at almost every balanced volume. If $h_\nu\le1$, then
\begin{equation}\label{eq:K-average}
  \int_{1/3}^{2/3}\mathfrak K_{\Sigma_p}\dd p\le C h_\nu^3 .
\end{equation}
\end{corollary}

\begin{proof}
For a concave profile with $I(0)=I(1)=0$, one has
$\sup_{[1/3,2/3]}I\le C h_\nu$ and
$\int_{1/3}^{2/3}(-I'')\dd p\le C h_\nu$ in the distributional sense. Combine these estimates with
\eqref{eq:profile-bound}.
\end{proof}

Thus the branch where $\mathfrak K_\Sigma$ is bounded below uniformly is not the generic branch in
a hypothetical small-Cheeger counterexample. The splitting branch is central, not exceptional.

\subsection{The weighted Reilly identity}

\begin{proposition}[Generalized Reilly identity]
\label{prop:reilly}
Let $\Omega\subset\R^n$ be a smooth bounded domain with boundary $\Sigma$, inner normal $n$,
and let $u$ be smooth on $\overline\Omega$. Then, with $L=\Delta-\nabla V\cdot\nabla$ and
$u_n=\partial_nu$,
\begin{equation}\label{eq:reilly}
  \int_\Omega(Lu)^2\dd\nu
  =\int_\Omega\Bigl(\norm{\nabla^2u}_\HS^2+\nabla u^T\nabla^2V\,\nabla u\Bigr)\dd\nu
  +\int_\Sigma\Bigl(H_\nu\,u_n^2+2u_n\,L_\Sigma u
  +II_\Sigma(\nabla_\Sigma u,\nabla_\Sigma u)\Bigr)\dd\sigma_\nu ,
\end{equation}
where $L_\Sigma$ is the induced weighted Laplacian on $\Sigma$. See
\cite{Milman2009Isoperimetric,Rosales2014} and the references therein for the weighted form.
In particular, for log-concave $\nu$ the interior terms are nonnegative and
\eqref{eq:reilly} converts boundary trace data of $\nabla u$ into the interior Dirichlet
quantity $\int(Lu)^2\dd\nu$, up to curvature boundary terms.
\end{proposition}

\begin{remark}[Dirichlet form of the trace estimate]
\label{rem:dirichlet-trace}
Apply \eqref{eq:reilly} on $\Omega=E$ to the Poisson solution $u_M$ of
Lemma~\ref{lem:boundary-rep}, with centered quadratic data $f_M$. The left side is computable:
$\int(Lu_M)^2\dd\nu=\Var_\nu^{E}$-type quantities of the quadratic form, of size controlled by
fourth moments, i.e.\ by thin-shell/quadratic-chaos information (Section~\ref{sec:qcts}). The
interior terms have a sign. The boundary terms involve $u_n^2$ weighted by $H_\nu$, mixed terms,
and the second fundamental form --- exactly the quantities that the stability inequality
\eqref{eq:stability} controls for near-minimizers, except on the constant mode. The Schur energy
\eqref{eq:Jsharp-def} is the bookkeeping device for this residual. The weighted Stein-trace
estimate is therefore reduced, at the level of mechanism, to: (a) quadratic-chaos control of the
interior Dirichlet quantity --- available up to the logarithmic ceiling of Section~\ref{sec:qcts};
(b) stability control of the mean-zero boundary modes --- available in model cases from
\eqref{eq:stability}; and (c) the constant boundary mode --- addressed either by
\eqref{eq:Jsharp-control} when $\mathfrak K_\Sigma$ is large or by splitting when it is small. We
emphasize that this is a reduction of mechanism, not yet a proof: the quantitative chaining of
(a)--(c) with universal constants, uniformly along the localization and at the covariance-weighted
scale, is open (Question~\ref{q:stein-weighted}).
\end{remark}

\subsection{Splitting: exact model statements}

\begin{proposition}[Exact splitting along a flat, potential-flat direction]
\label{prop:exact-splitting}
Let $\nu=e^{-V}dx$ be log-concave with convex support $K$ and $\theta\in S^{n-1}$, and suppose
$K$ is a cylinder in the $\theta$-direction: in coordinates $x=(y,z)$, $z=\inner x\theta$,
$K=K_1\times J$ with $K_1\subset\theta^\perp$ convex and $J\subseteq\R$ an interval. If
$\nabla^2V(\theta,\cdot)\equiv0$ on $\operatorname{int}K$, then the potential splits as
$V(y,z)=V_1(y)+cz$ for some $c\in\R$, and $\nu=\nu_1\otimes\nu_2$ is a product of a log-concave
measure $\nu_1\propto e^{-V_1}$ on $K_1$ and a log-affine factor $\nu_2\propto e^{-cz}$ on $J$;
normalizability of $\nu_2$ forces $J$ to be a proper interval --- a half-line when $c\ne0$, or a
bounded interval --- and not all of $\R$. Conversely, if $E=\{z\le z_0\}$ is a halfspace cut
orthogonal to such a direction, then $\Sigma$ is totally geodesic ($II\equiv0$),
$\nabla^2V(n,n)\equiv0$ on $\Sigma$, and the constant-mode curvature vanishes:
$\mathfrak K_\Sigma=0$.
\end{proposition}

\begin{proof}
On $\operatorname{int}K$, $\nabla^2V(\theta,\cdot)=0$ means $\partial_z\nabla V\equiv0$, so every
$\partial_{x_j}V$ is independent of $z$; in particular $\partial_zV$ is independent of $z$, and
since $\partial_{x_j}(\partial_zV)=0$ it is independent of $y$ as well, hence a constant $c$.
Integrating on the cylinder $K_1\times J$ gives $V=V_1(y)+cz$, so $\nu=\nu_1\otimes\nu_2$ with
$\nu_2\propto e^{-cz}$ on $J$, a finite measure only if $J\ne\R$ (a half-line when $c\ne0$, or a
bounded interval). The converse statements are immediate from $II=0$ for a hyperplane and
$\qJac_\Sigma=\abs{II}^2+\nabla^2V(n,n)=0$ on $\Sigma$.
\end{proof}

\begin{proposition}[Persistence of splitting under localization]
\label{prop:persistent-splitting}
If $V$ splits across $\theta$ in the sense of
Proposition~\textup{\ref{prop:exact-splitting}} (additively, $V(y,z)=V_1(y)+V_2(z)$ with
$V_2$ convex), then so does the localized potential at every time:
$V_t(x)=V(x)+\tfrac t2\abs x^2-c_t\cdot x$ splits as
$\bigl(V_1(y)+\tfrac t2\abs y^2-c_t'\cdot y\bigr)+\bigl(V_2(z)+\tfrac t2z^2-c_t''z\bigr)$.
In particular, the splitting direction is preserved pathwise, the covariance $A_t$ is
block-diagonal across $\theta^\perp\oplus\R\theta$ for all $t$, and halfspace cuts orthogonal
to $\theta$ have $\delta_t\parallel\theta$ and rank-one $K_t$ for all $t$.
\end{proposition}

\begin{proof}
The tilt $\tfrac t2\abs x^2-c_t\cdot x$ is itself additively separable in any orthogonal
decomposition. The block structure of $A_t$ and the alignment of $\delta_t$ follow from the
product structure of $\mu_t$ and of the cut.
\end{proof}

\begin{remark}[Boundary Obata and quantitative splitting]
\label{rem:obata}
Propositions~\ref{prop:exact-splitting}--\ref{prop:persistent-splitting} identify the equality
structure of the constant mode: $\mathfrak K_\Sigma=0$ for a minimizing halfspace cut forces,
along $\Sigma$, a totally geodesic boundary and a normal direction in which the potential has
no convexity --- the infinitesimal signature of a product with a log-affine factor. The
splitting philosophy of the geometric route is the corresponding stability assertion: a
near-worst measure with a near-minimal cut of small constant-mode curvature should be
quantitatively close, along the relevant directions, to such a split structure --- where the
two-color quantities are rank-one, the per-direction Carleson estimate
(Corollary~\ref{cor:per-direction}) already controls the source, and the variance dynamics are
one-dimensional. A rigidity statement of Obata type (vanishing constant-mode curvature implies
exact splitting, under appropriate completeness and minimality hypotheses) is the model case;
the quantitative version, with constants uniform along the localization, is
Question~\ref{q:splitting}. We do not claim any quantitative splitting theorem here.
\end{remark}
\end{document}
